\section{Discussion}

\subsection{Synthesis and Novelty}

In the introduction we presented prior work and how it informed GIAANN, focusing on background and synthesis.
Here, in the discussion we critically analyse how GIAANN diverges from those priors and what unique contributions it makes.

The GIAANN prototype brings together all the above themes into a unified architecture, and in doing so, it differs from prior works in significant ways while also addressing some of their limitations.
A key distinguishing feature is GIAANN’s explicit combination of multiple brain-inspired mechanisms in one system.
Where most previous models or theories focused on one or two of these aspects in isolation (e.g., HTM on columns and sequences, or dendritic ANNs on neuron substructures, or Hopfield nets on associative memory), GIAANN integrates semantic graphs, Hopfield dynamics, sequential spiking, dendritic compartments, concept cells, and columnar organisation into a single algorithmic framework.
This integrative approach is itself novel---it treats these concepts not as disparate research threads but as synergistic components of general intelligence.

Concretely, compared to its direct precursor, HFNLP \citep{bairesearch2024HFNLPpy}, GIAANN introduces Concept Columns.
HFNLP demonstrated that a Hopfield-like network could autonomously predict the next element in a sequence (particularly in language processing) via associative recall, but it treated the neuron population as largely homogeneous.
There was no structural separation for different types of information---all ``concepts'' lived in one recurrent pool, meaning the network struggled with interference and scaling (adding more concepts would eventually confuse the associative memory).
GIAANN overcomes this by dividing the neural substrate into columns, each handling a subset of the knowledge domain.
This means that when predicting the next neuron, GIAANN does not search the entire space blindly; it knows, for instance, that if the last active neuron was in a ``subject'' column, the next likely neuron might be in a related ``object'' column.
This structured traversal is far more scalable and mirrors how language or reasoning often has slots or roles (subject, predicate, object; or context, entity, attribute) that constrain what comes next.
In essence, GIAANN adds a contextual guide to the Hopfield recall: the semantic network (relational graph) plus column structure narrows the attractor search to sensible regions, thus solving one shortcoming of pure Hopfield nets, which can get spurious memories or require full connectivity.
The result is more reliable and interpretable sequence generation.

Relative to HTM and other cortical column models, GIAANN is less biologically literal in some ways (HTM models detailed minicolumns, layers, etc.), but it goes beyond them by incorporating full semantics and learning of explicit concepts.
HTM, for example, did not assign semantic labels to columns; it learned generic spatial-temporal patterns.
GIAANN builds in a semantic network from the start, ensuring each column/neuron has a meaning in human-understandable terms (at least after training).
This marriage of symbolic semantic structures (graphs) with sub-symbolic neural dynamics is reminiscent of neurosymbolic AI, but GIAANN does it in a tightly integrated, brain-inspired manner rather than attaching a logic engine onto a neural net.
The result is a system that can store knowledge (relations, facts) and process (chain through them) coherently.
Moreover, GIAANN’s use of dendritic logic within neurons is something HTM’s neuron model (though more complex than a perceptron) did not explore in depth.
Thus, GIAANN’s neurons are more powerful, potentially reducing the number of neurons needed and enabling the network to handle more complexity with sparse activity.

Against the backdrop of modern deep learning (e.g. Transformers), GIAANN stands out by eschewing backpropagation and dense continuous representations.
Instead, it relies on local learning rules (Hebbian) and highly sparse, event-driven computation.
This means GIAANN might not reach the performance of a Transformer on benchmark datasets and evaluation metrics, however, it offers advantages in interpretability, incremental learning, and energy efficiency (and experiments are indicative of consistent scaling).
Every activation in GIAANN can be traced to a concept, every connection to a relation---a transparency that is largely missing in black-box deep nets.
And whereas Transformers require enormous training data and then struggle with one-shot or continual learning, GIAANN (like the brain) can potentially incorporate new knowledge by adjusting a few connections in the semantic graph or adding a new neuron/column, without retraining the whole network.
This addresses the plasticity-stability dilemma in a structural way: knowledge is compartmentalised, so new learning in one column minimally interferes with established memories in others.

\subsection{Limitations}

It is important to note that while GIAANN draws on biology, it is not a direct simulation of the brain but rather a brain-inspired algorithm.
Consequently, some differences exist: for instance, real cortical columns may not each hold a distinct ontology of a concept---the brain’s columns may be more sensorily grounded whereas GIAANN’s may be more abstract and conceptual.
Also, inhibitory interneurons are critical in cortex for shaping activity; GIAANN currently emphasises excitatory drives with simplified global inhibition or as a mechanism for restricting prediction targets (exclusive top-1 selection, akin to winner-take-all \citep{maass2000winnerTakeAll}).
Additionally, GIAANN’s reliance on a predefined semantic cues (POS identification for concept column delimitation) could be seen as a limitation compared to deep learning models that learn representations purely from data.
However, one might argue GIAANN is simply incorporating prior knowledge (thing/relation distinction) which is a strength, not a weakness, especially for AGI aspirations.

\subsection{Conclusion}

In summary, GIAANN represents a novel synthesis that addresses shortcomings in prior art: it tackles the scaling issue of Hopfield memories by structuring them into columns; it improves interpretability and incremental learning relative to end-to-end deep nets; it increases neuron computational capacity instead of sheer quantity, aligning with energy constraints; and it grounds neural processing in explicit semantics, narrowing the gap between sub-symbolic neural nets and symbolic reasoning.
Its design is unorthodox in the contemporary ML landscape, but it hews closely to many cognitive and neuro principles, potentially offering a path to more general, human-like intelligence in machines.
As such, GIAANN differs from prior models: it is neither purely a Hopfield net, nor an HTM, nor a standard ANN, but an amalgam that leverages their strengths while mitigating their weaknesses.
The result is an architecture well-suited to embody understanding, reasoning, and learning in a way that is both computationally efficient and interpretable---a promising development in the quest for general AI.